

\subsection{Action Detection}
\label{sec:action_detection}
% \wsq{人体动作行为识别与时序定位}

Action detection (AD) finds human actions in untrimmed video: for each action instance it returns the time interval $[t_s, t_e]$ in which it occurs and its category $c$. Where action recognition classifies a clip that has already been trimmed to one action, detection must also decide when actions start and end, in streams that contain many of them.
It belongs to spatial structure construction because it organizes the entities and relations of a scene along the time axis into events, the temporal counterpart of the scene graphs above and the basis of video SGG and human--object interaction.
Figure~\ref{fig:p3_s6_3_4_ad_overview} illustrates the task.


\begin{figure}[!ht]
    \centering
    \includegraphics[width=\linewidth]{figures/part3/section6/action_detection.pdf}
    \caption{Action detection: each action instance in an untrimmed video is localized to $[t_s,t_e]$ and labeled with class $c$ (top). \textit{Bottom}: three task illustrations: (i) action classification of a trimmed clip, (ii) temporal action localization (illustrated with one interval), and (iii) temporal action detection of multiple instances.}
    \label{fig:p3_s6_3_4_ad_overview}
\end{figure}

\subsubsection{Task Definition}
\label{sec:6_3_4_definition}
Given an untrimmed video (or a continuous multi-sensor stream), AD outputs a set of tuples $(t_s, t_e, c)$ with start time, end time, and category for every action instance.
The following terms emphasize different aspects of the task. \textbf{Action classification (AC)} takes a trimmed clip and predicts one category, with no localization. \textbf{Temporal action localization (TAL)} predicts start and end times in untrimmed video, and may include multiple instances; the term is often used interchangeably with TAD. \textbf{Temporal action detection (TAD)} emphasizes joint localization and classification of multiple instances per video, possibly overlapping or concurrent, across temporal scales, each to be localized and classified.


\begin{figure}[!ht]
    \centering
    \includegraphics[width=0.99\linewidth]{figures/part3/section6/action_detection_method_1.pdf}
    \caption{Representative methods and mechanisms for (i) action classification, (ii) temporal localization with boundary estimation and alternative supervision signals, and (iii) temporal detection with anchor-based, anchor-free, boundary-based, and query-based formulations. Method timelines show publication order, not computational dependencies; spacing is not proportional. Same-year methods are grouped. The localization and detection panels emphasize different aspects of closely related tasks.}
    \label{fig:p3_s6_3_4_methodology_overview}
\end{figure}

\subsubsection{Methodology}

Figure~\ref{fig:p3_s6_3_4_methodology_overview} organizes the methods by the three settings; each builds on the previous one's representations.

\paragraph{Action Classification.}
Two-stream CNNs~\cite{simonyan2014two} model appearance from RGB frames in one stream and motion from dense optical flow in another, fusing the two at the output. 3D CNNs~\cite{diba2017temporal} convolve over space and time jointly, C3D~\cite{tran2015learning} and I3D~\cite{carreira2017quo} being the standard examples, at the cost of a large parameter count that R(2+1)D~\cite{tran2018closer} and P3D~\cite{qiu2017learning} reduce by factoring 3D convolution into 2D spatial and 1D temporal steps. TSN~\cite{wang2016temporal} samples a few snippets sparsely across a long video and supervises at video level, which scales long-range temporal modeling cheaply, and Transformers~\cite{girdhar2019video,yang2022recurring} model global spatio-temporal dependencies among sampled tokens beyond the local reach of convolution.

\paragraph{Temporal Action Localization.}
Early TAL ran a classifier over dense, fixed-size sliding windows, which was both slow and blind to the wide range of action durations. Structured segment networks~\cite{zhao2017temporal} instead model a proposal's starting, ongoing, and ending stages and score its completeness, which refines boundary regression. Boundary-sensitive methods~\cite{lin2018bsn} go further and predict start and end probabilities at every snippet, then pair high-confidence boundary points into proposals; this flexible boundary generation is the component on which later detection systems rest.

Because precise boundary annotation is expensive, weakly supervised TAL (WSTAL)~\cite{shou2018autoloc,zhang2021cola,paul2018w} learns from video-level labels alone. The task becomes multiple-instance learning: a video is positive if at least one segment contains the action, class activation sequences are derived from video features, and foreground segments are thresholded into proposals. Trained only with a classification loss, such models fire on the most discriminative moments rather than the whole action, so proposals are incomplete and high confidence does not imply accurate boundaries. Point-level supervision (P-TAL)~\cite{xia2024realigning,lee2021learning,fu2022compact} adds one labeled frame per instance, which tends to fall near the temporally most salient, roughly central moment, and uses it to realign confidence and infer complete boundaries from that single anchor.


\paragraph{Temporal Action Detection.}
TAD methods differ in how they generate and score temporal proposals.
Anchor-based detection~\cite{gao2017turn} transplants Faster R-CNN and SSD to the time axis: a dense set of multi-scale temporal anchors is classified and regressed, first in two stages (S-CNN~\cite{shou2016temporal}) and then in one (SSAD~\cite{lin2017single}); fixed anchors produce many redundant proposals and limit boundary precision, which anchor-free variants such as AFSD~\cite{lin2021learning} address by regressing boundaries from a temporal center representation.
Boundary-based detection~\cite{xiong2017pursuit,lin2018bsn} is bottom-up: start and end probabilities are predicted over the whole timeline and boundary matching scores every candidate pairing, which gives flexible, precise boundaries; TSA-Net~\cite{gong2019scale} adds multi-dilation temporal convolution for multi-scale context and BGCN~\cite{chen2021boundary} graph convolution over co-occurring proposals.
Query-based detection~\cite{wu2021towards,shi2022react,kim2023self,zhu2024dual} follows DETR: a small set of learned queries attends to global video features in a Transformer encoder--decoder and decodes the action set in parallel, removing anchors, NMS, and hand-designed matching at the cost of heavier computation.
Across the three, anchors give dense parallel scoring but coarse boundaries, boundary methods give precise boundaries but complex matching, and query methods give the simplest end-to-end pipeline; all three keep moving toward longer-range context, continuous temporal representations, and less frame-level annotation.

\subsubsection{Datasets and Metrics}
Table~\ref{tab:action_detection_dataset} lists the benchmarks by role. UCF101~\cite{soomro2012ucf101} and HMDB51~\cite{kuehne2011hmdb} are trimmed-clip classification sets; AVA~\cite{gu2018ava} adds spatio-temporal boxes on movie footage. THUMOS'14~\cite{jiang2014thumos}, whose 20 classes carry dense temporal boundaries, and ActivityNet~\cite{caba2015activitynet}, with 200 classes over 20K untrimmed videos, are the standard TAL and TAD benchmarks, extended in scale and granularity by HACS Segment~\cite{zhao2019hacs} and FineAction~\cite{liu2022fineaction}. MultiTHUMOS~\cite{yeung2018every}, Charades~\cite{sigurdsson2016hollywood}, and TVSeries~\cite{de2016online} supply dense, multi-label, overlapping actions in sports, homes, and television; Breakfast~\cite{kuehne2014language}, 50Salads~\cite{stein2013combining}, MPII Cooking 2~\cite{rohrbach2016recognizing}, and COIN~\cite{tang2019coin} cover procedural cooking and instructional activity for segmentation and step detection; and MEVA~\cite{corona2021meva} (thousands of hours of surveillance) and MoVi~\cite{ghorbani2021movi} (motion-captured actors) sit at the large-scale and controlled extremes.

\begin{table}[!t]
\centering
\caption{Benchmark datasets for action detection.}
\label{tab:action_detection_dataset}
\fontsize{9}{11}\selectfont
\setlength{\tabcolsep}{5.5pt}
\begin{tabular}{lccccc}
\toprule
\multirow{2}{*}{\textbf{Dataset}} & \textbf{Action} & \multirow{2}{*}{\textbf{Videos}} & \textbf{Duration} & \textbf{Multi-label} & \multirow{2}{*}{\textbf{Domain}} \\
& \textbf{Categories} & & \textbf{(hours)} & \textbf{(Per-Frame)} & \\
\midrule
UCF101~\cite{soomro2012ucf101}           & 101  & 13,320 & 27 & --  & Sports, Daily Activities\\
HMDB51~\cite{kuehne2011hmdb}           & 51   & 6,849  & -- & --  & Movies \\
AVA~\cite{gu2018ava}              & 80   & 430    & 108 & Yes & Movies \\
THUMOS'14~\cite{jiang2014thumos}        & 101 (20 w/ boundaries)  & 18,394 & --  & No & Sports \\
MultiTHUMOS~\cite{yeung2018every}  & 65                         & 413             & --                       & Yes         & Sports                 \\
Breakfast~\cite{kuehne2014language}    & 48                         & 1,712           & 77                    & No          & Cooking                \\
50Salads~\cite{stein2013combining}     & 17                         & 50              & 5.3                       & No          & Cooking                \\
MPII cooking 2~\cite{rohrbach2016recognizing} & 59                        & 273             & 27          & No          & Cooking                \\
TVSeries~\cite{de2016online}     & 30                         & 27              & --                   & Yes         & TV Series              \\
HACS Segment~\cite{zhao2019hacs} & 200                        & 50,000             & --                       & No          & Daily Activities       \\
ActivityNet-1.3~\cite{caba2015activitynet}  & 200  & 19,994 & 849   & No & Daily Activities \\
Charades~\cite{sigurdsson2016hollywood}         & 157  & 9,848  & --    & Yes         & Daily Activities \\
MoVi~\cite{ghorbani2021movi}             & 20   & --     & 9    & -- & Daily Activities \\
MEVA~\cite{corona2021meva}             & 37   & --     & 9,300  & -- & Surveillance \\
COIN~\cite{tang2019coin}         & 180 tasks                  & 11,827          & 476                       & No          & Daily Activities       \\
FineAction~\cite{liu2022fineaction}      & 106  & 16,732 & --    & Yes         & Daily Activities \\
\bottomrule
\end{tabular}
\end{table}

Each setting has its own metric. Action classification reports Top-1 accuracy, $\mathrm{Acc}_{\mathrm{Top}\text{-}1} = \frac{1}{N} \sum_{i=1}^{N} \mathbb{I}(\hat{y}_i = y_i)$, with Top-$k$ on large label sets; temporal boundaries are given, so localization is not assessed. Localization quality is measured by the temporal intersection over union of predicted and ground-truth intervals,
\begin{equation}
    \mathrm{tIoU}(I_p, I_g) =
    \frac{|I_p \cap I_g|}{|I_p \cup I_g|},
\end{equation}
a prediction counting as correct above a threshold $\tau$; TAL reports average precision at one or several thresholds. TAD scores detections $(I_p, \hat{y}, s)$ with confidence $s$ by mean average precision over classes and thresholds,
\begin{equation}
    \mathrm{mAP} = \frac{1}{|\mathcal{C}|}
    \sum_{c \in \mathcal{C}}
    \frac{1}{|\{\tau_m\}|}
    \sum_{\tau_m} \mathrm{AP}_c(\tau_m),
\end{equation}
which penalizes misclassification, loose boundaries, and duplicates together; THUMOS'14 uses $\tau \in \{0.3,\dots,0.7\}$ and ActivityNet $\{0.5,\dots,0.95\}$, and proposal methods also report average recall at a fixed proposal budget. Weakly supervised methods are often evaluated on ActivityNet-1.2 rather than 1.3, so their numbers are not interchangeable with fully supervised ones.

\begin{table*}[!t]
    \centering
    \caption{Representative methods under the two shared protocols. \textit{Top}: action classification, Top-1 accuracy (\%) on UCF101 (split 1, two-stream RGB+flow where applicable, with ImageNet or Kinetics pretraining as in each paper). \textit{Bottom}: temporal action detection on THUMOS'14, mAP (\%) at tIoU $=0.5$ and averaged over tIoU $\in\{0.3,\dots,0.7\}$, as reported by the cited papers and the comparisons in AFSD~\cite{lin2021learning}, CoLA~\cite{zhang2021cola}, and DualDETR~\cite{zhu2024dual}; ``weak'' = video-level labels only. ``--'' = not reported.}
    \label{tab:ad_performance}
    \fontsize{9}{11}\selectfont
    \setlength{\tabcolsep}{2.0mm}
    \begin{tabular}{L{0.22\textwidth} L{0.30\textwidth} c C{0.14\textwidth} C{0.16\textwidth}}
    \toprule
    \textbf{Method} & \textbf{Family} & \textbf{Year} & \textbf{UCF101 Top-1}$\uparrow$ & \\
    \midrule
    Two-Stream CNN~\cite{simonyan2014two} & RGB + optical-flow streams & 2014 & 88.0 & \\
    TSN~\cite{wang2016temporal}           & Sparse segment sampling    & 2016 & 94.2 & \\
    R(2+1)D~\cite{tran2018closer}         & Factorized 3D convolution  & 2018 & 97.3 & \\
    I3D~\cite{carreira2017quo}            & Inflated 3D convolution    & 2017 & \textbf{98.0} & \\
    \midrule
    \hiderowcolors \textbf{Method} & \textbf{Family} & \textbf{Year} & \textbf{mAP@0.5}$\uparrow$ & \textbf{Avg.\ mAP@[0.3:0.7]}$\uparrow$ \\
    \midrule
    S-CNN~\cite{shou2016temporal}     & Anchor-based, two-stage          & 2016 & 19.0 & -- \\
    SSN~\cite{zhao2017temporal}       & Segment-structure proposals      & 2017 & 29.8 & -- \\
    BSN~\cite{lin2018bsn}             & Boundary-sensitive proposals     & 2018 & 36.9 & 36.8 \\
    AutoLoc~\cite{shou2018autoloc}    & Weak, outer--inner contrast      & 2018 & 21.2 & -- \\
    CoLA~\cite{zhang2021cola}         & Weak, snippet contrastive        & 2021 & 32.2 & -- \\
    AFSD~\cite{lin2021learning}       & Anchor-free, salient boundary    & 2021 & 55.5 & 52.0 \\
    DualDETR~\cite{zhu2024dual}       & Query-based (dual DETR)          & 2024 & \textbf{70.4} & \textbf{66.8} \\
    \bottomrule
    \end{tabular}
\end{table*}

Table~\ref{tab:ad_performance} shows how far the two settings have diverged. Trimmed-clip classification on UCF101 was effectively solved once Kinetics-pretrained 3D networks arrived, with I3D and R(2+1)D within a point of each other near 98\%. Detection on THUMOS'14 tells a different story: mAP@0.5 moved from 19.0 (S-CNN) to 36.9 (BSN) over the anchor-to-boundary transition, jumped to 55.5 with anchor-free salient-boundary features (AFSD), and reached 70.4 with query-based decoding (DualDETR), while the best weakly supervised method (CoLA, 32.2) matches fully supervised accuracy of only a few years earlier, which quantifies what video-level labels still cost.


\subsubsection{Challenges}
Fully supervised TAD depends on densely annotated boundaries, which scale poorly to long, loosely structured real-world video; weakly and semi-supervised methods relax this with video-level labels or sparse temporal cues~\cite{zhao2025pose,ghoddoosian2023weakly}, at the accuracy cost the table above makes visible.
Actions are also still detected in isolation. Coupling TAD with structured temporal representations such as scene graphs or event schemas~\cite{wu2025universal,lu2025benchmarking} would let a system infer boundaries, temporal relations, and higher-level activity structure jointly, which matters most in instructional video, driving, and human--robot interaction.
Finally, benchmarks and detectors remain visual and closed-set, whereas real action understanding draws on audio, language, gaze, and context and must handle actions outside a fixed vocabulary~\cite{zhou2024learning,bao2025exploiting,liu2026synergistic}; multimodal temporal detection, where audio or language sharpens boundaries and disambiguates semantics, and open-vocabulary formulations that loosen the classification-centric setup are the directions large multimodal video corpora now make feasible.
